
\textbf{Instance-Level Detection.} Carlini et al. (2021) and Ye et al. (2022) demonstrated membership inference on memorized sequences. Fu et al. (2024) established fundamental limitations: MIA achieves AUC≈50\% in pretraining because models learn distributions. Semantic similarity detection fails against adversarial paraphrasing: Dekoninck et al. (2024) showed EAL achieves TPR<2\% at 1\% FPR for semantic detectors while preserving ~15\% gains.

\textbf{Learning Dynamics.} Swayamdipta et al. (2020) categorized training examples by confidence dynamics. Toneva et al. (2019) identified unforgettable examples. Feldman (2020) connected memorization to atypical examples. Our TSG probe approach was conceptually grounded in differential learning rates, but this mechanism was never implemented.

\textbf{Loss Landscape Geometry.} Li et al. (2018) visualized loss surfaces. Fort and Jastrzebski (2019) connected Hessian spectrum to generalization. Our Tier 3 geometric detection was conceptually grounded in this literature but was never implemented.

\textbf{Continuous Evaluation.} Jain et al. (2024) introduced LiveCodeBench with 400+ programming problems weekly. Kiela et al. (2021) proposed Dynabench. These approaches prevent contamination through constant benchmark refreshment. Our approach attempted complementary training-time detection but failed before testing.

\textbf{Position.} We attempted to extend detection beyond instance-level features toward learning trajectory analysis. The novelty was Task Signature Graphs as paraphrase-invariant representations and geometric inevitability—the hypothesis that gains require detectable alignment. However, implementation failure prevented testing. Prior work (Fu et al., Dekoninck et al.) remains state-of-the-art.
